\chapter{Progettazione logica}

\section{Stima del volume dei dati}

Si fornisce in questa fase una tabella contenente il numero medio di istanze di ogni entità e associazione dello schema concettuale finale.

\begin{table}[H]
\centering
\renewcommand{\arraystretch}{1.3} % Spaziatura comoda per le righe
\setlength{\tabcolsep}{10pt}     % Spaziatura orizzontale nelle celle

\begin{tabularx}{\textwidth}{|X|c|r|}
\hline
\textbf{Concetto} & \textbf{Tipo (E/R)} & \textbf{Volume} \\
\hline
UTENTE & E & 10.000 \\
\hline
pubblicazione & R & 150.000 \\
\hline
RICETTA & E & 150.000 \\
\hline
composizione & R & 1.500.000 \\
\hline
INGREDIENTE & E & 1.000 \\
\hline
appartenenza & R & 450.000 \\
\hline
CATEGORIA & E & 30 \\
\hline
richiesta & R & 300.000 \\
\hline
ORDINE & E & 300.000 \\
\hline
acquisto & R & 600.000 \\
\hline
PROMOZIONE & E & 1.000 \\
\hline
offerta & R & 15.000 \\
\hline
SCONTO & E & 15.000 \\
\hline
applicazione & R & 15.000 \\
\hline
recensione & R & 3.000.000 \\
\hline
\end{tabularx}
\caption*{\textbf{Tabella 3.1.1} - \textit{Tavola dei volumi del modello E/R.}}
\end{table}


\section{Descrizione delle operazioni principali e stima della loro frequenza}

In questa fase si propone una tavola delle operazioni utilizzata per costituire una stima delle principali
operazioni richieste da utenti standard e amministratori.

\begin{table}[H]
\centering
\renewcommand{\arraystretch}{1.3}
\setlength{\tabcolsep}{5pt} % Riduce leggermente lo spazio tra le colonne per far stare tutto

\small % Riduce leggermente la dimensione del testo per un layout più pulito

% Usiamo due colonne 'X' per bilanciare lo spazio tra Nome Operazione e Frequenza
\begin{tabularx}{\textwidth}{|c|>{\raggedright\arraybackslash}X|>{\raggedright\arraybackslash}X|c|}
\hline
\textbf{Codice Op.} & \textbf{Nome operazione} & \textbf{Frequenza} & \textbf{Tipo} \\
\hline
U1 & REGISTRAZIONE UTENTE & 5 al giorno & I \\
\hline
U2 & VISUALIZZAZIONE RICETTE & 300 al giorno & I \\
\hline
U3 & VISUALIZZAZIONE RICETTE PER FILTRO & (1000 utenti attivi) $\times$ 1 = 1000 al giorno & I \\
\hline
U4 & VISUALIZZAZIONE INGREDIENTI & 300 al giorno & I \\
\hline
U5 & PUBBLICAZIONE RECENSIONE & (1000 utenti attivi) $\times$ 0.1 = 100 al giorno & I \\
\hline
U6 & PUBBLICAZIONE RICETTA & 50 al giorno & I \\
\hline
U7 & VISUALIZZAZIONE VANTAGGI & (1000 utenti attivi) $\times$ 0.5 = 500 al giorno & I \\
\hline
U8.1 & VISUALIZZAZIONE MIGLIORI RICETTE IN BASE ALLE RECENSIONI & (1000 utenti attivi) $\times$ 1.5 = 1500 al giorno & I \\
\hline
U8.2 & VISUALIZZAZIONE MIGLIORI UTENTI & 10 al giorno & I \\
\hline
U8.3 & VISUALIZZAZIONE RICETTE PIÙ ORDINATE & (1000 utenti attivi) $\times$ 0.1 = 100 al giorno & I \\
\hline
U8.4 & VISUALIZZAZIONE CATEGORIE DI RICETTA PIÙ ORDINATE & (1000 utenti attivi) $\times$ 0.1 = 100 al giorno & I \\
\hline
U9 & EFFETTUAZIONE ORDINE & 100 al giorno & I \\
\hline
U10 & VISUALIZZAZIONE STORICO ORDINI & 50 al giorno & I \\
\hline
U11 & VISUALIZZAZIONE SPECIFICA RICETTA & (1000 utenti attivi) $\times$ 1.5 = 1500 al giorno & I \\
\hline
A1 & INSERIMENTO INGREDIENTE & 1 al mese & I \\
\hline
A2 & AGGIORNAMENTO INGREDIENTE & 5 a settimana & I \\
\hline
A3 & INSERIMENTO CATEGORIA & 1 all’anno & I \\
\hline
A4 & LANCIO PROMOZIONE & 5 all’anno & I \\
\hline
A5 & PUBBLICAZIONE SCONTO & 75 all’anno & I \\
\hline
A6 & RIMOZIONE UTENTI & 1 al mese & B \\
\hline
\end{tabularx}
\caption*{\textbf{Tabella 3.2.1} - \textit{Tavola delle operazioni del sistema.}}
\end{table}

\textbf{Nota}: la frequenza indicata nell’operazione U3 indica la frequenza totale di tutte le operazioni di ricerca per filtro (U3.1 - U3.6), le cui singole frequenze sono approssimativamente in ugual parti.